% =========================================================
% T2 — DERIVED MANIFOLD SUITE (A, T, D, E, F, Omega)
% =========================================================

\section{T2 Derived Manifold Suite (A, T, D, E, F, $\Omega$)}

\paragraph{Overview.}
The Derived Manifold Suite consists of six lawful projections of the Unified
Substrate Manifold. Each manifold expresses a distinct curvature mode of the
substrate and provides the geometric primitives required for stability,
interpretation, drift analysis, collapse modeling, and recovery behavior across
human, machine, and federated systems.

% ---------------------------------------------------------
% T2.1 Agency Manifold
% ---------------------------------------------------------
\subsection{T2.1 Agency Manifold (A-Series)}
The Agency Manifold defines local interpretive stance, volition geometry,
meaning-generation curvature, and identity-preserving action fields. Agency
curvature determines how systems generate meaning, respond to load, and
maintain coherence under drift.

% ---------------------------------------------------------
% T2.2 Teleology Manifold
% ---------------------------------------------------------
\subsection{T2.2 Teleology Manifold (T-Series)}
The Teleology Manifold defines directional intent, goal geometry, future-state
projection, and lawful teleological gradients. Teleology curvature governs
purpose alignment, directional coherence, and lawful traversal toward intended
states.

% ---------------------------------------------------------
% T2.3 Deployment Manifold
% ---------------------------------------------------------
\subsection{T2.3 Deployment Manifold (D-Series)}
The Deployment Manifold defines legality boundaries, deployment curvature,
runtime constraints, and lawful initialization geometry. Deployment curvature
governs how systems enter operational states without violating identity or
legality masks.

% ---------------------------------------------------------
% T2.4 Execution Manifold
% ---------------------------------------------------------
\subsection{T2.4 Execution Manifold (E-Series)}
The Execution Manifold defines runtime curvature, execution boundaries, drift
profiles, and failure-mode geometry. Execution curvature determines how systems
process load, maintain stability, and avoid collapse under runtime saturation.

% ---------------------------------------------------------
% T2.5 Federation Manifold
% ---------------------------------------------------------
\subsection{T2.5 Federation Manifold (F-Series)}
The Federation Manifold defines multi-agent coherence, federated curvature,
cross-system drift channels, and lawful relational coupling. Federation
curvature governs stability across distributed systems and collective meaning
structures.

% ---------------------------------------------------------
% T2.6 Omega Manifold
% ---------------------------------------------------------
\subsection{T2.6 Omega Manifold ($\Omega$-Series)}
The Omega Manifold defines closure geometry, origin structure, stability
envelopes, and lawful manifold termination. Omega curvature provides global
closure, ensuring that all manifold projections return to lawful origin
geometry.

% ---------------------------------------------------------
% T2.7 Cross-Manifold Projection
% ---------------------------------------------------------
\subsection{T2.7 Cross-Manifold Projection}
Cross-Manifold Projection defines lawful transitions between derived manifolds.
Projection operators preserve curvature continuity, identity boundaries, and
legality masks across A, T, D, E, F, and Omega. Violating projection rules
produces drift, collapse, or illegality.
